\section{Extending MOSFETs with Shift Register}

\subsection{Overview}

The Raspberry Pi Zero provides a limited number of GPIO pins. Where multiple independent
loads must be switched, a serial-to-parallel shift register allows eight MOSFET gates to
be controlled from three GPIO pins. The device chosen is the \textbf{74HC595}, an
8-bit serial-in, parallel-out shift register that operates from \SI{3.3}{\volt} and is
widely available in an SMD package.

Data is clocked into the shift register serially and latched to the eight parallel outputs
simultaneously, ensuring that all outputs update at the same instant. Although the
74HC595 can be driven by bit-banging three GPIO pins, the implementation here uses the
hardware \textbf{SPI1} peripheral on the Raspberry Pi Zero for reliable timing and
minimal CPU overhead. Only one MOSFET is energised at a time; the controlling byte
therefore always has exactly one bit set.

\subsection{Circuit Description}

The Raspberry Pi Zero SPI1 peripheral provides three signals: MOSI (data), SCLK (clock),
and CE (chip enable, used here as the latch signal). Each signal passes through a
\SI{470}{\ohm} series resistor before reaching the 74HC595 to protect the GPIO pins
from transient currents.

The 74HC595 output enable pin (\textoverline{OE}) is tied permanently to ground so
that the parallel outputs are always active. The master reset pin (\textoverline{MR})
is tied permanently to \SI{3.3}{\volt} to prevent accidental clearing of the register.
A \SI{100}{\nano\farad} bypass capacitor is placed directly across the Vcc and GND
pins of the IC.

Each of the eight parallel outputs drives a MOSFET gate through a \SI{500}{\ohm}
series resistor. A \SI{10}{\kilo\ohm} resistor from gate to ground ensures the gate
is pulled firmly low when the corresponding output is de-asserted, preventing the
MOSFET from floating into an undefined state.

\subsection{Circuit Diagram}

\begin{verbatim}
  Raspberry Pi Zero                  74HC595 (SOIC-16)
  (SPI1)                             +----------------+
                          3.3V -----| Vcc(16) (8)GND |------ GND
                                    |                |
   GPIO 20 (MOSI) --[ 470R ]--------| DS   (pin 14)  |   [SER]
   GPIO 21 (SCLK) --[ 470R ]--------| SHCP (pin 11)  |   [SRCLK]
   GPIO 18 (CE0)  --[ 470R ]--------| STCP (pin 12)  |   [RCLK]
                                    |                |
                            GND ----| /OE  (pin 13)  |  (always enabled)
                           3.3V ----| /MR  (pin 10)  |  (master reset off)[SRCLR]
                                    |                |
                                    |   Q0  (pin 15) |--+
                                    |   Q1  (pin 1)  |--+  (Q1..Q7 identical)
                                    |   Q2  (pin 2)  |--+
                                    |   Q3  (pin 3)  |--+
                                    |   Q4  (pin 4)  |--+
                                    |   Q5  (pin 5)  |--+
                                    |   Q6  (pin 6)  |--+
                                    |   Q7  (pin 7)  |--+
                                    |       (pin 9)  |--x  (Q7S serial data out)
                          100nF     |                |
                                    +----------------+
  --||-- Vcc to GND, close to IC
 
  One MOSFET channel (repeated x8):

  74HC595 Qn ---[ 500R ]---+--- MOSFET Gate (AO3400A or similar)
                            |
                          10k
                            |
                           GND                MOSFET Source --- GND
                                              MOSFET Drain  --- Load (-)
\end{verbatim}

\subsection{Parts List}

\begin{table}[h]
\centering
\begin{tabular}{llll}
\hline
\textbf{Reference} & \textbf{Part Number} & \textbf{Description} & \textbf{Package} \\
\hline
U1      & 74HC595           & 8-bit shift register, serial-in parallel-out  & SOIC-16 \\
Q1--Q8  & AO3400A           & N-channel MOSFET, \SI{30}{\volt}, \SI{5.7}{\ampere} & SOT-23 \\
R1--R3  & R470-0603         & \SI{470}{\ohm} GPIO series resistors          & 0603 \\
R4--R11 & R500-0603         & \SI{500}{\ohm} MOSFET gate resistors          & 0603 \\
R12--R19 & R10K-0603        & \SI{10}{\kilo\ohm} gate pull-down resistors   & 0603 \\
C1      & 100nF ceramic     & Bypass capacitor, Vcc to GND                  & 0603 \\
\hline
\end{tabular}
\caption{Shift register MOSFET controller components}
\label{tab:shiftreg_parts}
\end{table}

\subsection{Programming Notes}

The SPI1 peripheral on the Raspberry Pi Zero is accessed via the \texttt{spidev} library.
SPI mode 0 (CPOL=0, CPHA=0) is compatible with the 74HC595. A clock speed of
\SI{500}{\kilo\hertz} is more than adequate; the 74HC595 is rated to several megahertz
but there is no benefit to running faster for low-speed MOSFET switching.

A single byte is transferred per operation. The byte has exactly one bit set to select
one MOSFET channel, or is zero to de-energise all channels. The \texttt{xfer2()} call
clocks out the data and automatically toggles the chip enable line, which is connected
to the latch (STCP) pin, causing the outputs to update at the end of the transfer.

Note that SPI1 on the Raspberry Pi Zero requires the overlay \texttt{spi1-1cs} to be
enabled in \texttt{/boot/config.txt}:

\begin{verbatim}
    dtoverlay=spi1-1cs
\end{verbatim}

\subsection{Example Code}

\begin{verbatim}
import spidev

SPI_BUS = 1       # SPI1
SPI_DEVICE = 0    # CE0

spi = spidev.SpiDev()
spi.open(SPI_BUS, SPI_DEVICE)
spi.max_speed_hz = 500000
spi.mode = 0


def set_mosfet(channel: int) -> None:
    """
    Energise a single MOSFET by channel number.

    Parameters
    ----------
    channel : int
        Channel to activate, 0 through 7.
        Pass -1 or any out-of-range value to de-energise all channels.
    """
    if 0 <= channel <= 7:
        byte = 1 << channel
    else:
        byte = 0
    spi.xfer2([byte])


def all_off() -> None:
    """De-energise all MOSFET channels."""
    spi.xfer2([0x00])
\end{verbatim}

\noindent To activate channel 4 then turn everything off:

\begin{verbatim}
    set_mosfet(4)
    time.sleep(1.0)
    all_off()
\end{verbatim}


% ---------------------------------------------------------------------------
\section{High-Side Control of 5\,V Load with MOSFET}

\subsection{Overview}

In circuits requiring both high-side and low-side switching, the load is interrupted at
both the positive supply terminal and the ground return. The low-side switch is an
N-channel MOSFET (AO3400A), which is straightforward to drive from a \SI{3.3}{\volt}
GPIO. The high-side switch requires a P-channel MOSFET because its source terminal
connects to the positive supply rail.

The P-channel MOSFET chosen is the \textbf{AO3401A}, the complementary P-channel
counterpart to the AO3400A. Both devices share the SOT-23 package and are produced
by the same manufacturer, making them a well-matched pair. The AO3401A is rated at
\SI{-30}{\volt} Vds and \SI{-4}{\ampere} continuous drain current, well in excess of
the anticipated load current.

\subsection{The Level-Shifting Problem}

A P-channel MOSFET turns on when its gate is pulled sufficiently below its source
voltage. With the source at \SI{+5}{\volt}, the gate must be pulled close to ground
to turn the device on, and must be pulled to \SI{+5}{\volt} to turn it off. A
\SI{3.3}{\volt} GPIO cannot swing to \SI{+5}{\volt} and therefore cannot turn the
device fully off. A small-signal NPN transistor is used to solve this: it pulls the
gate to ground when the GPIO is high (turning the P-channel on), and allows the gate
to be pulled to \SI{+5}{\volt} through a resistor when the GPIO is low (turning the
P-channel off).

\subsection{Circuit Description}

The GPIO drives the base of an MMBT3904 NPN transistor through a \SI{10}{\kilo\ohm}
base resistor. The collector of the NPN connects to the gate of the AO3401A. A
\SI{10}{\kilo\ohm} resistor from the gate to \SI{+5}{\volt} pulls the gate high
when the NPN is off, ensuring the P-channel is firmly off. When the NPN saturates,
it pulls the gate to ground, turning the P-channel on.

The logic is inverted by the NPN stage: GPIO high turns the high-side MOSFET on,
and GPIO low turns it off. This matches the behaviour of the low-side N-channel
switch, so both MOSFETs are enabled by a high logic level on their respective GPIO
pins. The load is energised only when both GPIO pins are high simultaneously.

When the supply voltage is \SI{5}{\volt} and the gate is pulled to ground by the NPN,
Vgs of the AO3401A is \SI{0}{\volt} $-$ \SI{5}{\volt} $=$ \SI{-5}{\volt}. The
AO3401A gate threshold is \SI{-1}{\volt}, so the device is driven well into
conduction.

\subsection{Circuit Diagram}

\begin{verbatim}
  High-side control:

  RPi GPIO A --[ 10k ]--+-- MMBT3904 Base
                         |
                    MMBT3904 Emitter ----------- GND
                    MMBT3904 Collector --+------ AO3401A Gate
                                        |
                                       10k
                                        |
                                      +5V

  +5V ----------- AO3401A Source
                  AO3401A Drain ----+---- Load (+)
                                    |
                                  Load (-)
                                    |
                               AO3400A Drain
                               AO3400A Source --- GND
                               AO3400A Gate --[ 500R ]--+-- RPi GPIO B
                                                        |
                                                       10k
                                                        |
                                                       GND

  Load energised when: GPIO A = HIGH and GPIO B = HIGH
\end{verbatim}

\subsection{Parts List}

\begin{table}[h]
\centering
\begin{tabular}{llll}
\hline
\textbf{Reference} & \textbf{Part Number} & \textbf{Description} & \textbf{Package} \\
\hline
Q1 & AO3401A   & P-channel MOSFET, \SI{-30}{\volt}, \SI{-4}{\ampere}, high-side switch & SOT-23 \\
Q2 & AO3400A   & N-channel MOSFET, \SI{30}{\volt}, \SI{5.7}{\ampere}, low-side switch  & SOT-23 \\
Q3 & MMBT3904  & NPN BJT, \SI{40}{\volt}, \SI{200}{\milli\ampere}, gate driver         & SOT-23 \\
R1 & R10K-0603 & \SI{10}{\kilo\ohm} NPN base resistor                                  & 0603   \\
R2 & R10K-0603 & \SI{10}{\kilo\ohm} P-channel gate pull-up to \SI{5}{\volt}            & 0603   \\
R3 & R500-0603 & \SI{500}{\ohm} N-channel gate series resistor                         & 0603   \\
R4 & R10K-0603 & \SI{10}{\kilo\ohm} N-channel gate pull-down to GND                   & 0603   \\
\hline
\end{tabular}
\caption{High-side and low-side MOSFET switch components}
\label{tab:highside_parts}
\end{table}

\subsection{Logic Table}

\begin{table}[h]
\centering
\begin{tabular}{ccccl}
\hline
\textbf{GPIO A} & \textbf{NPN} & \textbf{AO3401A} & \textbf{GPIO B (AO3400A)} & \textbf{Load} \\
(high-side)     &              & (P-channel)      & (low-side)                &               \\
\hline
LOW  & OFF & OFF & LOW  & Off \\
LOW  & OFF & OFF & HIGH & Off \\
HIGH & ON  & ON  & LOW  & Off \\
HIGH & ON  & ON  & HIGH & \textbf{On} \\
\hline
\end{tabular}
\caption{High-side and low-side switch logic}
\label{tab:highside_logic}
\end{table}

\subsection{Example Code}

\begin{verbatim}
import RPi.GPIO as GPIO

HIGH_SIDE_PIN = 23   # controls AO3401A via MMBT3904
LOW_SIDE_PIN  = 24   # controls AO3400A directly

GPIO.setmode(GPIO.BCM)
GPIO.setup(HIGH_SIDE_PIN, GPIO.OUT, initial=GPIO.LOW)
GPIO.setup(LOW_SIDE_PIN,  GPIO.OUT, initial=GPIO.LOW)


def load_on() -> None:
    """Energise the load by enabling both high-side and low-side switches."""
    GPIO.output(HIGH_SIDE_PIN, GPIO.HIGH)
    GPIO.output(LOW_SIDE_PIN,  GPIO.HIGH)


def load_off() -> None:
    """De-energise the load. Low side is opened first."""
    GPIO.output(LOW_SIDE_PIN,  GPIO.LOW)
    GPIO.output(HIGH_SIDE_PIN, GPIO.LOW)
\end{verbatim}

\noindent The low-side switch is opened first in \texttt{load\_off()} to interrupt
current flow before the high side is released. This order avoids a brief condition
where the high side is off but the low side is still on, which would momentarily
short the load's negative terminal through the N-channel with no current path,
producing no harm but making the sequence logically cleaner.
